\section{Progettazione}
\label{sec:progettazione}

\subsection{Architettura generale}

Il sistema è composto da due parti distinte: l'applicazione Android, che funge da client, e un backend REST containerizzato, che gestisce autenticazione, sincronizzazione della history e la logica dei gruppi.

\subsubsection*{Architettura Android -- pattern MVVM}

Il lato Android adotta il pattern \textbf{MVVM} (Model--View--ViewModel). I tre livelli principali sono:

\begin{itemize}[label={--}, leftmargin=1.2cm]
    \item \textbf{UI layer} (\texttt{ui/}): ogni schermata è un \emph{Composable} che osserva uno \texttt{StateFlow} esposto dal proprio ViewModel. Quando lo stato cambia, Compose effettua una recomposition, cambiando la UI in base al suo stato. Il color scheme e la tipografia generale dell'intera UI sono definiti da un tema \textbf{Material 3}.

    \item \textbf{ViewModel} (\texttt{ui/<schermata>/}): raccoglie i dati dalle repository, applica la logica di presentazione e li espone tramite una \texttt{data class} di stato (\texttt{UiState}). Un esempio si può osservare in \texttt{HomeViewModel}, dove vengono esposte le ultime 3 entry della history:

\inputcodeblock{home_viewmodel_recent.kt}

    \item \textbf{Data layer} (\texttt{data/}): espone i dati tramite interfacce di repository, nascondendo alla UI se i dati provengono dal database locale o dal backend.
\end{itemize}

\subsubsection*{Data layer -- struttura a strati}

Il data layer è organizzato in tre sotto-moduli:

\begin{itemize}[label={--}, leftmargin=1.2cm]
    \item \texttt{data/local/db/}: database SQLite creato con Room per la persistenza locale dei dati. E' il metodo primario di persistenza dell'applicazione. 
    \item \texttt{data/remote/}: interfaccia Retrofit \texttt{LocalShareApi} con le chiamate REST per la comunicazione con il web service Flask, il quale persiste i dati su un database non relazionale MongoDB. 
    \item \texttt{data/p2p/}: gestione della connessione Wi-Fi Direct e del trasferimento file con protocollo TCP.
\end{itemize}

\subsubsection*{Comunicazione con il backend}

La sincronizzazione con il backend è \textbf{opzionale}: l'app funziona in modo autonomo anche senza connessione, poiché l'identità del dispositivo è un UUID generato localmente alla prima installazione tramite \texttt{DeviceIdProvider}, senza bisogno di contattare alcun server. Il backend viene coinvolto solo per la sincronizzazione della history e la condivisione tramite gruppi.

Quando la sincronizzazione è abilitata, le chiamate REST sono effettuate tramite \textbf{Retrofit}, con \textbf{kotlinx.serialization} come converter per il parsing JSON.
L'autenticazione avviene tramite JWT: l'app registra il proprio UUID sul backend e ottiene un token, che viene salvato nelle \texttt{Shared\-Preferences}; per far sì che ogni chiamata abbia nell'header il token, viene utilizzato un interceptor:

\inputcodeblock{okhttp_jwt_interceptor.kt}

La repository \texttt{SyncedHistoryRepository} si occupa della sincronizzazione tra la repository locale e quella backend, dando precedenza alla persistenza locale: ogni trasferimento viene prima salvato nel database Room, poi si tenta l'invio al backend. Se la rete non è disponibile, un \texttt{WorkManager} si occuperà di ritentantare la sincronizzazione una volta che il dispositivo ritornerà connesso.

\subsection{Struttura dei dati / Database}

L'applicazione utilizza due livelli di persistenza distinti: un database locale sul dispositivo Android (metodo principale e sempre utilizzato) e un database remoto sul backend (secondario e utilizzato secondo le preferenze dell'utente).

\subsubsection*{Persistenza locale -- Room (SQLite)}

Il database locale necessita di una sola tabella contenente le entry che costituiscono la history delle transazioni di un utente: \texttt{history}, mappata dalla \texttt{data class} \texttt{HistoryEntity}:

\begin{itemize}[label={--}, leftmargin=1.2cm]
    \item \texttt{transferId}: UUID che identifica univocamente il trasferimento; è la chiave condivisa con il backend.
    \item \texttt{peerDeviceId}: UUID del dispositivo con cui è avvenuto lo scambio.
    \item \texttt{fileName}: nome del file trasferito.
    \item \texttt{sizeBytes}: dimensione del file trasferito.
    \item \texttt{direction}: enum \texttt{SENT} / \texttt{RECEIVED} per identificare i file trasferiti da quelli ricevuti.
    \item \texttt{timestampMillis}: timestamp della transazione.
    \item \texttt{latitude}, \texttt{longitude}: campi opzionali che salvano la posizione in cui il trasfermento è avvenuto.
    \item \texttt{isSynced}: flag presente solamente nel DB locale che indica se l'entry è già stata inviata al backend.
\end{itemize}

\subsubsection*{Persistenza remota -- MongoDB}

Nel backend sono presenti 3 collection MongoDB:

\begin{itemize}[label={--}, leftmargin=1.2cm]
    \item \texttt{users}: un documento per dispositivo registrato, contenente lo \texttt{user\_id} (UUID) e il timestamp \texttt{created\_at}.

    \item \texttt{transfers}: un documento per trasferimento, ha gli stessi campi di \texttt{HistoryEntity} locale (escluso \texttt{isSynced}) più \texttt{user\_id} per indicarne il proprietario e \texttt{group\_code} opzionale.

    \item \texttt{groups}: un documento per gruppo, con un codice alfanumerico di 6 caratteri generato casualmente e un array \texttt{members} contenente gli \texttt{user\_id} dei partecipanti. Una volta entrato in un gruppo, un dispositivo vedrà nella sua sezione history le transazioni degli altri dispositivi del gruppo. 
\end{itemize}

\subsection{Scelte progettuali}

\subsubsection*{Design pattern}

\begin{itemize}[label={--}, leftmargin=1.2cm]
    \item \textbf{Repository pattern}: Ogni repository (locale e backend) è nascosta grazie ad un'interfaccia: \texttt{DeviceRepository} e \texttt{HistoryRepository}. Grazie a questo layer di astrazione, i viewModel istanziano l'interfaccia e quindi non sanno i dettagli del database su cui stanno lavorando.  

    \item \textbf{Strategy pattern}: \texttt{P2pManager} e \texttt{TransferManager} sono interfacce con implementazioni concrete (\texttt{WifiDirectP2pManager}, \texttt{SocketTransferManager}) iniettate da \texttt{AppContainer}. Questo layer di astrazione nasconde l'implementazione della metodologia di trasferimento dati dalla logica di business. Se per esempio si volesse implementare un nuovo metodo di trasferimento (es. Bluetooth)
    per applicarla alla logica di business basterebbe cambiare l'injection in \texttt{AppContainer}. 

    \item \textbf{Singleton}: Nel progetto si trovano due \texttt{object} Kotlin. \texttt{TransferSessionState} espone un flag condiviso \texttt{isTransferring} usato per coordinare \texttt{Incoming\-Transfer\-Watcher} e \texttt{Transfer\-Foreground\-Service} (il suo ruolo è approfondito nella sezione \hyperref[subsec:broadcast]{BroadcastReceiver}). \texttt{AppContainer} è anch'esso un singleton, descritto successivamente.

    \item \textbf{Factory Method}: ogni schermata definisce una \texttt{viewModelFactory} che incapsula la logica di costruzione del proprio ViewModel, separandola dal codice che lo utilizza. Questo è necessario perché Compose non può costruire ViewModel con parametri arbitrari tramite il meccanismo di default; la factory fornisce al sistema le istruzioni su come creare l'istanza. Un esempio concreto in \texttt{HomeScreen}:

\begin{codeblock}[language=Kotlin]
private val HomeViewModelFactory = viewModelFactory {
    initializer {
        HomeViewModel(
            AppContainer.deviceIdProvider,
            AppContainer.deviceRepository,
            AppContainer.historyRepository,
        )
    }
}
\end{codeblock}

    \item \textbf{Dependency injection manuale}: \texttt{AppContainer} funge da injector, istanziando tutte le dipendenze principali (repository, client HTTP, database, manager P2P) e le fornisce ai ViewModel tramite le factory descritte sopra.
\end{itemize}

\subsubsection*{Composable}

Un \textbf{Composable} è una funzione Kotlin annotata con \texttt{@Composable} che descrive una porzione dell'interfaccia utente in modo dichiarativo: invece che essere modifiata esplicitamente, la view intercetta e cambia a seconda dello stato corrente attraverso una recomposition.

Lo stato può essere gestito a due livelli distinti:

Lo \textbf{stato di ViewModel} (\texttt{StateFlow}) è usato per dati che devono sopravvivere alle recomposition, alle rotazioni di schermo e che provengono dalla logica di business.

Lo \textbf{stato interno al Composable} viene gestito con \texttt{remember \{ mutableStateOf(...) \}}. E' usato per dati esclusivamente visivi di cui non è richiesta la sopravvivenza alla recomposition. Un esempio concreto all'interno del progetto è il seguente:

\begin{itemize}[label={--}, leftmargin=1.2cm]
    \item le \texttt{HistoryCard} all'interno della schermata di History usano \texttt{var expanded by remember \{ mutableStateOf(false) \}} per tracciare il loro stato, che può essere espanso o collassato.
\end{itemize}

\begin{figure}[h]
    \centering
    \begin{minipage}[c]{0.42\textwidth}
        \centering
        \includegraphics[width=\textwidth]{card_history_collapsed.png}
    \end{minipage}
    \hfill
    \begin{minipage}[c]{0.42\textwidth}
        \centering
        \includegraphics[width=\textwidth]{card_history_expanded.png}
    \end{minipage}
    \caption{HistoryCard collassata (sinistra) ed espansa (destra)}
\end{figure}

\subsubsection*{Coroutines}

Le operazioni asincrone (chiamate di rete, I/O su disco, trasferimento file) sono gestite tramite coroutine.
Nei ViewModel le coroutine vengono lanciate in \texttt{viewModelScope}, che le annulla automaticamente quando il ViewModel viene distrutto.\\
Un utilizzo più avanzato delle coroutine si ha nella funzione \texttt{observePeers} utilizzata al momento della discovery dei dispositivi nelle vicinanze dove l'API di WiFi-Direct, basata su callback, viene adattata a un \texttt{Flow} Kotlin:

\begin{codeblock}[language=Kotlin]
override fun observePeers(): Flow<List<PeerDevice>> = callbackFlow {
    val receiver = object : BroadcastReceiver() {
        override fun onReceive(ctx: Context, intent: Intent) {
            if (intent.action == WifiP2pManager.WIFI_P2P_PEERS_CHANGED_ACTION)
                requestPeers()
        }
    }
    ContextCompat.registerReceiver(context, receiver, IntentFilter(...), ...)
    awaitClose { context.unregisterReceiver(receiver) }
}
\end{codeblock}

\subsubsection*{BroadcastReceiver e Wi-Fi Direct}%
\label{subsec:broadcast}

L'API Wi-Fi Direct notifica gli eventi di rete tramite \textbf{Broad\-cast\-Receiver}. In \texttt{Wifi\-Direct\-P2p\-Manager} vengono registrati due receiver:

\begin{itemize}[label={--}, leftmargin=1.2cm]
    \item \texttt{WIFI\_P2P\_PEERS\_CHANGED\_ACTION}: inviato dal sistema quando la lista dei dispositivi vicini cambia; alla ricezione viene richiesta la lista aggiornata e inviata al Flow.
    \item \texttt{WIFI\_P2P\_CONNECTION\_CHANGED\_ACTION}: inviato quando lo stato della connessione P2P cambia; alla ricezione vengono richieste le informazioni sul gruppo formato (indirizzo del group owner, ruolo del dispositivo).
\end{itemize}

\texttt{Incoming\-Transfer\-Watcher} osserva il Flow di \texttt{observeConnectionInfo()} in una coroutine con scope dedicato e, quando rileva una nuova connessione P2P, avvia il \texttt{Transfer\-Foreground\-Service} in modalità ricezione. Tuttavia il broadcast \texttt{WIFI\_P2P\_CONNECTION\_CHANGED\_ACTION} viene ricevuto da entrambi i dispositivi coinvolti nella connessione, incluso il mittente. Per evitare che il mittente avvii erroneamente una sessione di ricezione, \texttt{Incoming\-Transfer\-Watcher} controlla il flag \texttt{Transfer\-Session\-State.isTransferring} prima di procedere: se il dispositivo ha già avviato un trasferimento, il flag è \texttt{true} e il servizio non viene avviato.

\subsubsection*{Foreground Service}

Il trasferimento file avviene al di fuori del ciclo di vita di Activity e ViewModel, in un \textbf{Foreground Service} (\texttt{TransferForegroundService}). Questa scelta è necessaria perché il trasferimento può durare diversi secondi e deve continuare anche se l'utente naviga fuori dalla schermata. Android richiede che un servizio in primo piano mostri una \textbf{notifica persistente} per tutta la durata dell'operazione, in modo che l'utente sia consapevole dell'attività in corso.
